\documentclass{resume}

\usepackage[left=0.4 in,top=0.4in,right=0.4 in,bottom=0.4in]{geometry}
\usepackage{ctex}
\usepackage{xcolor}
\usepackage{fontspec}
\defaultfontfeatures{Extension = .otf}
\usepackage{fontawesome}
\newcommand{\tab}[1]{\hspace{.2667\textwidth}\rlap{#1}} 
\newcommand{\itab}[1]{\hspace{0em}\rlap{#1}}

% \hypersetup{%
%   colorlinks=false,% hyperlinks will be black
%   linkbordercolor=red,% hyperlink borders will be red
%   pdfborderstyle={/S/U/W 1}% border style will be underline of width 1pt
% }

\name{Tzu Chun, Hsu}
\address{New Taipei City, Taiwan} 
\address{
\small \href{tel:+886-987605719}{ \raisebox{-0.1\height}\faPhone\ \underline{+886-987605719} ~} \href{mailto:vm3y3rmp40719@gmail.com}{\raisebox{-0.2\height}\faEnvelope\  \underline{vm3y3rmp40719@gmail.com}} ~ 
  \href{https://www.linkedin.com/in/tzu-chun-hsu-ab4b3b188/}{\raisebox{-0.2\height}\faLinkedinSquare\ \underline{tzu-chun-hsu-ab4b3b188}}  ~
  \href{https://github.com/Oscarshu0719}{\raisebox{-0.2\height}\faGithub\ \underline{Oscarshu0719}}
}

\begin{document}

%----------------------------------------------------------------------------------------
% WORK EXPERIENCE
%----------------------------------------------------------------------------------------
\vspace{-1em}
\begin{rSection}{Work Experience}
    {\bf Camera ISP Firmware Engineer, Realtek Semiconductor Corp.} \hfill {02 2025 -- present}\\
    {\textit{AI camera algorithm development}} \hfill \textit{Hsinchu, Taiwan}\\
    $\cdot$ Developed high-resolution image stitching solution for smart glasses to ensure consistent ISP pipeline behavior between streaming and snapshot modes.\\
    $\cdot$ Implemented ISP pipeline features for smart glasses with V4L2 and interrupt controls and optimized SDK flow for maintainability and efficiency.\\
    $\cdot$ Migrated ISP engine for smart glasses from Linux to FreeRTOS and implemented custom porting layer for unsupported interfaces.\\
    $\cdot$ Refactored core CModel architecture and consolidated sub-applications for improved maintainability and reduced code redundancy.\\
    $\cdot$ Implemented low-level control APIs with thread and cross-process synchronization for Windows for cross-platform compatibility.\\
    $\cdot$ Designed and optimized automated verification systems for ISP/IPU (FPGA/ASIC) hardware platforms, involving Python server and C agent on hardware platforms.\\
    $\cdot$ Migrated communication protocol from JRPC to protobuf for the Linux-to-FreeRTOS migration for cross-platform compatibility and improved transmission efficiency.
\end{rSection}

%----------------------------------------------------------------------------------------
% EDUCATION
%----------------------------------------------------------------------------------------
\vspace{-1em}
\begin{rSection}{Education}

{\bf National Yang Ming Chiao Tung University} \hfill {09 2022 -- 02 2025}\\
{\textit{Master of Science in Computer Science and Engineering}} \hfill \textit{Hsinchu, Taiwan}\\
$\cdot$ Overall GPA: 4.19/4.30

{\bf Zhejiang University} \hfill {09 2016 -- 07 2020}\\
{\textit{Bachelor of Engineering in Computer Science and Technology}} \hfill \textit{Hangzhou, China}\\
$\cdot$ Last two years GPA: 3.30/4.00

\end{rSection}

%----------------------------------------------------------------------------------------
% AWARDS
%----------------------------------------------------------------------------------------
\vspace{-1em}
\begin{rSection}{AWARDS}
    {\bf CCCC-Mobile Application Innovation Contest} \hfill {09 2018}\\
    {\textit{First Prize and Most Innovative Award}} \hfill \textit{Hangzhou, China}

    {\bf The 4th ”Internet+” Innovation and Entrepreneurship Competition} \hfill {09 2018}\\
    {\textit{Gold Award}} \hfill \textit{Hangzhou, China}

    {\bf College Students’ Innovative Entrepreneurial Training Plan Program (SRTP)} \hfill {04 2018 -- 05 2019}\\
    {\textit{Excellent}} \hfill \textit{Hangzhou, China}
\end{rSection} 

%----------------------------------------------------------------------------------------
% SKILLS	
%----------------------------------------------------------------------------------------
\vspace{-1em}
\begin{rSection}{SKILLS}
\begin{tabular}{ @{} >{\bfseries}l @{\hspace{6ex}} l }
    Languages & Python, C, C++, Java, JavaScript\\
    Technologies/Frameworks & Flask, PyTorch, Spring Boot, Vue.js
\end{tabular}\\
\end{rSection}

\newpage

%----------------------------------------------------------------------------------------
% PROJECTS
%----------------------------------------------------------------------------------------
\vspace{-1em}
\begin{rSection}{PROJECTS}
\vspace{-1.25em}
    \item \textbf{\href{https://docs.google.com/presentation/d/1797idRRmgeD2JnUslGbflmGUO2X5lMY6/edit?usp=drive_link&ouid=101248488395326982475&rtpof=true&sd=true}{\textbf{\large{\underline{Chord learning and adversarial framework for symbolic music generation}}}}} \hfill {04 2018 -- 05 2019}\\
    \mbox{} \hfill \textit{Hangzhou, China}

    This project developed an automated arrangement system for ancient zither music using a hybrid deep learning architecture. 
    We proposed a framework combining VAE (Variational Autoencoder) for structural feature extraction and GAN (Generative Adversarial Network) for enhancing creative variance. 
    To stabilize the training process and resolve issues like mode collapse, we introduced WGAN-GP and designed three customized loss functions. 
    Furthermore, We developed an algorithmic Zither Music Theory Filter to incorporate traditional pentatonic scales and performance techniques (e.g., tremolo, glissando) into the generative process. 
    The project resulted in a functional iOS App with chord visualization and was recognized with the First Prize and Most Innovative Award in the CCCC-Mobile Application Innovation Contest and 
    the Gold Award in the 4th "Internet+" Innovation and Entrepreneurship Competition.

    \item \textbf{\href{https://drive.google.com/file/d/1736XHtaeT58FL_gbOa9XjlyA6xAXDkg2/view?usp=drive_link}{\textbf{\large{\underline{Voice Conversion Based on Generative Adversarial Networks}}}}} \hfill {03 2020 -- 05 2020}\\
    \mbox{} \hfill \textit{Hangzhou, China}

    This project proposed an enhanced non-parallel voice conversion (VC) framework based on StarGAN-VC2 to improve the naturalness and quality of converted speech. 
    Key optimizations include transitioning acoustic features from traditional MCEPs to Mel-spectrograms and 
    replacing Conditional Instance Normalization (CIN) with Adaptive Instance Normalization (AdaIN) for more effective style transfer. 
    Additionally, WaveGlow, a flow-based neural vocoder, was integrated to replace the WORLD vocoder, significantly enhancing synthesized voice clarity. 
    Experimental results, validated by Mel-cepstral Distortion (MCD) and Mean Opinion Score (MOS), demonstrated that the proposed model outperforms the StarGAN-VC2 baseline in both objective acoustic similarity and subjective naturalness.
    
    \item \textbf{\href{https://drive.google.com/file/d/176xCWH6j0VAuoXLtB_HppxyZN-6ezq9_/view?usp=drive_link}{\textbf{\large{\underline{Synthetic Data Generation using Conditional Normalizing Flows}}}}} \hfill {03 2023 -- 05 2023}\\
    \mbox{} \hfill \textit{Hsinchu, Taiwan}

    This project introduced SurrogateFlow, a generative model based on Conditional Normalizing Flows (CNF) designed to create high-fidelity synthetic datasets. 
    Unlike GANs, SurrogateFlow leveraged the exact density estimation of CNFs to avoid mode collapse, ensuring the generated data maintains the full distribution and integrity of the original dataset. 
    A key innovation was the integration of anomaly scores into the training process, allowing users to adjust sampling probabilities to highlight outliers for scientific analysis. 
    Experimental results demonstrated that SurrogateFlow outperforms Gaussian Mixture Models (GMM) and GANs in terms of data generation quality, comprehensiveness, and the ability to handle privacy-sensitive data.

    \item \textbf{Floating image quality compensation algorithm technology} \hfill {08 2023 -- 02 2025}\\
    \mbox{} \hfill \textit{Hsinchu, Taiwan}

    This project addressed image degradation in light field floating 3D displays caused by hardware imperfections, such as misalignment, distortion, and ghosting. 
    I developed a two-stage compensation algorithm: (1) using Operator Splitting Quadratic Program (OSQP) to model the physical imaging process and pre-calculate lens intensity weights, and 
    (2) an iterative optimization framework to refine light field images based on a custom objective function. 
    By decoupling device-specific optical weights from image refinement, the proposed method significantly reduced computational overhead. 
    Experimental results demonstrated that this approach effectively compensates for manufacturing defects and display panel degradation, 
    aligning perceived image quality with theoretical optical standards.
\end{rSection}

\newpage

%----------------------------------------------------------------------------------------
% COURSES	
%----------------------------------------------------------------------------------------
\vspace{-1em}
\begin{rSection}{COURSES}
\subsection*{Undergraduate Course List:}

\noindent
\begin{minipage}[t]{0.31\textwidth}
    \subsubsection*{專業必修與核心課程}
    \begin{enumerate}
        \item 程序設計基礎
        \item 計算機科學基礎
        \item 離散數學及其應用
        \item 程序設計專題
        \item 數字邏輯設計
        \item 數據結構基礎
        \item 面向對象程序設計
        \item 面向信息技術的溝通技巧
        \item 數據庫系統
        \item 高級數據結構與算法分析
        \item 密碼學
        \item 計算機組成
        \item 操作系統
        \item 計算機體系結構
        \item 計算機網絡
        \item Java應用技術
        \item 計算理論
        \item 編譯原理
        \item 計算機科學思想史
        \item 軟件工程
        \item 專題研討
        \item 人工智能
        \item 博奕論原理與應用
        \item 智能終端軟件開發
        \item 應用運籌學基礎
        \item 畢業論文（設計）
    \end{enumerate}
\end{minipage}\hfill
\begin{minipage}[t]{0.31\textwidth}
    \subsubsection*{公共基礎課程}
    \begin{enumerate}
        \item 大學英語III
        \item 微積分（甲）I
        \item 線性代數（甲）
        \item 大學物理（甲）I
        \item 微積分（甲）II
        \item 大學物理實驗
        \item 大學物理（乙）II
        \item 概率論與數理統計
        \item 大學英語II
    \end{enumerate}
\end{minipage}\hfill
\begin{minipage}[t]{0.31\textwidth}
    \subsubsection*{通識課程}
    \begin{enumerate}
        \item 電工電子學概論
        \item 生物工程導論
        \item 攝影藝術
        \item 科學與生活
        \item 綠色化學與可持續發展
        \item Linux應用技術基礎
        \item Java程序設計
        \item 西方音樂簡史
        \item 生命科學
        \item 材料科學與人類文明
        \item 應用運籌學
        \item 服務科學導論
        \item 成長中的親密關係
        \item Python程序設計
    \end{enumerate}
\end{minipage}

\newpage

\vspace{-1em}
\subsection*{Master Course List:}

\noindent
\begin{enumerate}
    \item 資料視覺化與視覺分析
    \item 演算法
    \item 圖形理論
    \item 影像處理
    \item 深度學習
    \item 深度學習實驗
    \item 雲原生軟體開發與最佳實踐
    \item 影像編修技術與特效合成
\end{enumerate}

\end{rSection} 

\end{document}
